\section{Conclusion}

By linking newly available 1940-1950 census data 
with War Relocation Authority records,
I constructed a measure of the predicted probability of internment 
and compared post-war outcomes for Japanese Americans 
against a control group of Chinese Americans.
I sought to answer a puzzle posed by the case of Japanese internment:
how could the forced internment have paradoxically 
led to economic opportunity for some?
My results agree with surprising finding that a higher probability of internment
is associated with higher individual earnings in 1950, after resettlement had occurred.
However, this aggregate effect masks a critical source of heterogeneity;
the gains were driven almost entirely by women.
While those likely to be interned were more prone to migrate to higher-income and urban counties,
this channel only partially explains the results.
This could suggest that internment fundamentally altered the labor market trajectories of women
in ways not captured by simple occupational changes.

These findings build on a growing literature
that re-evaluates the long-run economic consequences of internment.
My work aligns with recent studies 
by \cite{arellano-boverDisplacementDiversityMobility2022} 
and \cite{shoagCausalEffectPlace2016} 
in documenting upward mobility,
but it adds a crucial dimension by revealing that these gains were not uniform.
The concentration of effects among women
suggests that the forced break from pre-war community structures
and labor market constraints on the West Coast
may have paradoxically opened new avenues for economic participation in the postwar economy.
Although internees lost their freedoms, property, and dignity,
these results provide a more complex understanding
of how individuals subjected to forced removal
can also adapt to their new environments.
